\documentclass{beamer}
\usetheme{Madrid}
\usecolortheme{default}

\usepackage[utf8]{inputenc}
\usepackage{listings}
\usepackage{xcolor}
\usepackage{tikz}
\usepackage{graphicx}

% Code listing style
\lstset{
    basicstyle=\ttfamily\footnotesize,
    keywordstyle=\color{blue}\bfseries,
    commentstyle=\color{green!60!black},
    stringstyle=\color{red},
    breaklines=true
}

\title{Introduction to Cryptography}
\subtitle{From Classical Ciphers to Modern Encryption}
\author{Computer Science Course}
\date{\today}

\begin{document}

\frame{\titlepage}

\begin{frame}{Table of Contents}
\tableofcontents
\end{frame}

% =====================================================
\section{Introduction to Cryptography}
% =====================================================

\begin{frame}{What is Cryptography?}
\begin{block}{Definition}
The science and practice of securing communication through codes and ciphers
\end{block}

\textbf{Etymology:} Greek \emph{kryptos} (hidden) + \emph{graphein} (to write) \pause

\vspace{0.5cm}
\textbf{We use cryptography every day:}
\begin{itemize}
    \item<2-> Messaging apps (WhatsApp, Signal)
    \item<3-> Online shopping
    \item<4-> Banking
    \item<5-> WiFi connections
\end{itemize}

% IMAGE SUGGESTION: Icon/diagram showing various everyday uses of cryptography
\end{frame}

\begin{frame}{Why Do We Need Cryptography?}
\textbf{Without cryptography, we are vulnerable to:}

\begin{description}
    \item[Eavesdropping] Unauthorized parties reading your messages
    \item[Tampering] Attackers modifying messages without detection
    \item[Impersonation] Someone pretending to be you
    \item[Data theft] Personal information being stolen\pause
\end{description}


\vspace{0.5cm}
\textbf{Cryptography provides:}
\begin{itemize}
    \item \textbf{Confidentiality:} Only authorized parties can read
    \item \textbf{Integrity:} Data cannot be modified undetected
    \item \textbf{Authentication:} Verify sender identity
    \item \textbf{Non-repudiation:} Sender cannot deny sending
\end{itemize}

% IMAGE SUGGESTION: Diagram showing attacker scenarios and cryptographic protections
\end{frame}

\begin{frame}{Cryptography's Limitations}


% IMAGE SUGGESTION: XKCD comic about voting software/security

\textbf{Even strong encryption cannot protect against:}
\begin{itemize}
    \item Social engineering attacks
    \item Malware on your device
    \item Poor password choices
    \item Insider threats
    \item Physical access to devices
    \item Implementation bugs
\end{itemize}\pause

\begin{alertblock}{Remember}
``Cryptography is like a really good lock on your door. It won't help if the walls are made of paper.''
\end{alertblock}
\end{frame}

\begin{frame}{Basic Terminology}
\begin{description}
    \item[Plaintext]<1-> Original, readable message
    \item[Ciphertext]<2-> Encrypted, unreadable message
    \item[Cipher]<3-> Algorithm for encryption/decryption
    \item[Key]<4-> Secret information used with cipher
    \item[Encryption]<5-> Converting plaintext to ciphertext
    \item[Decryption]<5-> Converting ciphertext to plaintext
    \item[Keyspace]<6-> Set of all possible keys
\end{description}

% IMAGE SUGGESTION: Simple diagram showing plaintext → encryption → ciphertext → decryption → plaintext
\end{frame}

% \begin{frame}{Meet Alice, Bob, and Eve}
% \textbf{Standard names in cryptography:}

% \begin{description}
%     \item[Alice] First party sending a message
%     \item[Bob] Second party receiving a message
%     \item[Eve] Eavesdropper trying to intercept
%     \item[Mallory] Malicious attacker modifying messages
%     \item[Trent] Trusted third party
% \end{description}

% % IMAGE SUGGESTION: Cartoon illustration of Alice, Bob, Eve showing communication flow
% \end{frame}

% =====================================================
\section{Classical Cryptography}
% =====================================================

\begin{frame}{Historical Context}
\textbf{Cryptography through history:}

\begin{itemize}
    \item \textbf{Ancient Egypt (1900 BCE):} Hieroglyphic substitutions
    \item \textbf{Spartans (5th century BCE):} Scytale transposition cipher
    \item \textbf{Julius Caesar (100-44 BCE):} Caesar cipher for military
    \item \textbf{Medieval Period:} Al-Kindi's frequency analysis (850 CE)
    \item \textbf{WWII:} Enigma machine and Bletchley Park
\end{itemize}

% IMAGE SUGGESTION: Timeline with historical cryptography examples/images
\end{frame}

\begin{frame}{Steganography vs. Cryptography}
\begin{block}{Steganography}
Hiding the \emph{existence} of a message
\end{block}

\textbf{Examples:}
\begin{itemize}
    \item Invisible ink
    \item Microdots
    \item Acrostics
    \item Data hidden in images/audio
\end{itemize}
\pause
\vspace{0.3cm}
\begin{exampleblock}{Try This}
\texttt{Since everyone can read, encoding text in neutral\\
sentences is doubtfully effective.}
\end{exampleblock}\pause

Hidden message: \textbf{SECRET INSIDE} (first letters)

% IMAGE SUGGESTION: Example of steganography - image with hidden data
\end{frame}

\begin{frame}[fragile]{Caesar Cipher}
\textbf{One of the oldest ciphers:}

Shift each letter by a fixed number of positions

\vspace{0.3cm}
\textbf{Example with shift k=3:}
\begin{verbatim}
Plaintext:  THE QUICK BROWN FOX
Ciphertext: WKH TXLFN EURZQ IRA
\end{verbatim}

\vspace{0.3cm}
\textbf{Mathematical notation:}
\begin{align*}
E(x) &= (x + k) \mod 26 \\
D(x) &= (x - k) \mod 26
\end{align*}

% IMAGE SUGGESTION: Caesar cipher wheel diagram
\end{frame}

\begin{frame}{Breaking the Caesar Cipher}
\textbf{Why Caesar cipher is weak:}

\begin{itemize}
    \item \textbf{Small keyspace:} Only 25 possible keys
    \item \textbf{Brute force trivial:} Try all shifts
    \item \textbf{Frequency analysis works:} Most common letter likely = E
\end{itemize}

\vspace{0.5cm}
\begin{exampleblock}{Try It Yourself}
Use \texttt{cryptii.com} to experiment with Caesar ciphers
\end{exampleblock}

% IMAGE SUGGESTION: Frequency analysis chart showing letter distribution
\end{frame}

\begin{frame}[fragile]{Monoalphabetic Substitution}
\textbf{Any permutation of alphabet:}

\begin{verbatim}
Plain:  ABCDEFGHIJKLMNOPQRSTUVWXYZ
Cipher: QWERTYUIOPASDFGHJKLZXCVBNM
\end{verbatim}

\textbf{Keyspace:} $26! \approx 4 \times 10^{26}$ (huge!)

\vspace{0.5cm}
\textbf{But still vulnerable to frequency analysis:}
\begin{itemize}
    \item Most common letters: E, T, A, O, I, N, S, H, R
    \item Common digrams: TH, HE, IN, ER, AN
    \item Common trigrams: THE, AND, ING
\end{itemize}

% IMAGE SUGGESTION: Frequency analysis bar chart for English text
\end{frame}

\begin{frame}[fragile]{Transposition Ciphers}
\textbf{Rearrange letters rather than replacing them}

\vspace{0.3cm}
\textbf{Rail Fence Cipher (3 rails):}

\begin{verbatim}
Message: WE ARE DISCOVERED FLEE AT ONCE

W . . . E . . . C . . . R . . . L . . . T . . . E
. E . R . D . S . O . E . E . F . E . A . O . C .
. . A . . . I . . . V . . . D . . . E . . . N . .

Ciphertext: WECRL TEERD SOEEF EAOCA IVDEN
\end{verbatim}

% IMAGE SUGGESTION: Visual diagram of rail fence zigzag pattern
\end{frame}

\begin{frame}[fragile]{Vigenère Cipher}
\textbf{Polyalphabetic substitution:} Different shifts at different positions

\vspace{0.3cm}
\textbf{Example with keyword LEMON:}
\begin{verbatim}
Plaintext:  A T T A C K A T D A W N
Keyword:    L E M O N L E M O N L E
Ciphertext: L X F O P V E F R N H R
\end{verbatim}

\vspace{0.3cm}
\begin{itemize}
    \item Called ``le chiffre indéchiffrable'' (the indecipherable cipher)
    \item Same plaintext letter → different ciphertext letters
    \item Defeats simple frequency analysis
\end{itemize}

% IMAGE SUGGESTION: Vigenère square (tabula recta)
\end{frame}

\begin{frame}{Breaking Vigenère}
\textbf{Despite being called ``indecipherable,'' it can be broken:}

\vspace{0.3cm}
\textbf{Method:}
\begin{enumerate}
    \item Find keyword length (Kasiski examination)
    \begin{itemize}
        \item Look for repeated sequences
        \item Distances between repetitions share factors with keyword length
    \end{itemize}
    \item Split ciphertext into groups (every $n$-th letter)
    \item Each group encrypted with same Caesar shift
    \item Use frequency analysis on each group
\end{enumerate}
\end{frame}

\begin{frame}[fragile]{One-Time Pad: Perfect Secrecy}
\textbf{The only provably unbreakable cipher!}

\vspace{0.3cm}
\textbf{Requirements:}
\begin{itemize}
    \item Key as long as message
    \item Truly random key
    \item Never reuse key (even partially)
    \item Destroy key after use
\end{itemize}
\pause
\vspace{0.3cm}
\begin{exampleblock}{Example}
\begin{verbatim}
Plaintext:  H E L L O
Key:        X M C K L (random, same length)
Ciphertext: E Q N V Z
\end{verbatim}
\end{exampleblock}

\textbf{Why unbreakable:} Any plaintext equally likely without key!
\end{frame}

\begin{frame}{Why We Don't Use One-Time Pad}
\textbf{Practical challenges:}

\begin{enumerate}
    \item \textbf{Key length:} Must equal message length (1 GB video = 1 GB key)
    \item \textbf{Key distribution:} How to securely share huge random keys?
    \item \textbf{True randomness:} Hard to generate
    \item \textbf{Key management:} Never reuse, must destroy
    \item \textbf{Synchronization:} Both parties must use same key portion
\end{enumerate}
\pause
\vspace{0.3cm}
\begin{alertblock}{Historical Failure}
Soviet Union reused one-time pads → NSA's VENONA project decrypted thousands of messages
\end{alertblock}
\end{frame}

% =====================================================
\section{Modern Cryptography}
% =====================================================

\begin{frame}{The Cryptographic Revolution}
\begin{block}{Kerckhoffs's Principle (1883)}
``A cryptosystem should be secure even if everything about the system, except the key, is public knowledge.''
\end{block}

\vspace{0.3cm}
\textbf{Modern approach benefits:}
\begin{itemize}
    \item Algorithms can be published and tested by experts
    \item Security relies only on secret keys
    \item Standard algorithms used worldwide
    \item If key compromised, just change key (not entire system)
\end{itemize}
\end{frame}

\begin{frame}{Classical vs. Modern Cryptography}
\begin{center}
\begin{tabular}{|l|l|l|}
\hline
\textbf{Aspect} & \textbf{Classical} & \textbf{Modern} \\
\hline
Security basis & Secret algorithm & Secret key \\
Key size & Small & Large (128-4096 bits) \\
Foundation & Simple math & Complex mathematics \\
Breakability & Pattern analysis & Computational hardness \\
Speed & Fast (by hand) & Requires computers \\
Standardization & Difficult & Widespread \\
\hline
\end{tabular}
\end{center}

% IMAGE SUGGESTION: Visual comparison diagram
\end{frame}

\begin{frame}{Computational Security}
\textbf{Modern cryptography relies on computational hardness}

\vspace{0.3cm}
\textbf{Example: 128-bit key}
\begin{itemize}
    \item $2^{128} \approx 3.4 \times 10^{38}$ possible values
    \item Testing 1 trillion keys/second
    \item Would take $\approx 10^{21}$ years
    \item \emph{Far longer than age of universe!}
\end{itemize}

\vspace{0.3cm}
\textbf{Key insight:}
Breaking is theoretically possible but practically infeasible
\end{frame}

\begin{frame}{Three Pillars of Modern Cryptography}
\begin{enumerate}
    \item<1-> \textbf{Symmetric Encryption}
    \begin{itemize}
        \item Same key for encryption/decryption
        \item Very fast
        \item Examples: AES, ChaCha20
    \end{itemize}

    \item<2-> \textbf{Asymmetric Encryption}
    \begin{itemize}
        \item Different keys (public/private)
        \item Solves key distribution
        \item Examples: RSA, ECC
    \end{itemize}

    \item<3-> \textbf{Cryptographic Hash Functions}
    \begin{itemize}
        \item One-way functions
        \item Fixed output size
        \item Examples: SHA-256, SHA-3
    \end{itemize}
\end{enumerate}

% IMAGE SUGGESTION: Diagram showing relationship between the three pillars
\end{frame}

% =====================================================
\section{Symmetric Encryption}
% =====================================================

\begin{frame}{Symmetric Encryption Overview}
\textbf{Same key for both operations}\pause

\vspace{0.3cm}
\textbf{Advantages:}
\begin{itemize}
    \item Very fast
    \item Strong security with short keys (256 bits)
    \item Low computational requirements
\end{itemize}
\pause
\vspace{0.3cm}
\textbf{Disadvantages:}
\begin{itemize}
    \item Key distribution problem
    \item Key management ($n$ people = $\frac{n(n-1)}{2}$ keys)
    \item Key compromise affects all messages
\end{itemize}

% IMAGE SUGGESTION: Diagram showing Alice and Bob with shared key
\end{frame}

\begin{frame}{Block Ciphers vs. Stream Ciphers}
\begin{columns}
\column{0.5\textwidth}
\textbf{Block Ciphers}
\begin{itemize}
    \item Fixed-size blocks (128 bits)
    \item Examples: AES, DES
    \item Need padding
    \item Require mode of operation
\end{itemize}
\pause

\column{0.5\textwidth}
\textbf{Stream Ciphers}
\begin{itemize}
    \item Bit-by-bit or byte-by-byte
    \item Examples: ChaCha20, RC4
    \item No padding needed
    \item Good for streaming data
\end{itemize}
\end{columns}

\vspace{0.5cm}
% IMAGE SUGGESTION: Visual comparison of block vs stream cipher operation
\end{frame}

\begin{frame}{Block Cipher Modes of Operation}
\textbf{Problem:} Identical plaintext blocks → identical ciphertext blocks

\vspace{0.3cm}
\textbf{Common modes:}
\begin{description}
    \item[ECB] Electronic Codebook - \alert{NOT SECURE!}
    \item[CBC] Cipher Block Chaining - each block XORed with previous
    \item[CTR] Counter Mode - turns block cipher into stream cipher
    \item[GCM] Galois/Counter Mode - encryption + authentication
\end{description}

% IMAGE SUGGESTION: ECB Penguin example showing why ECB is insecure
\end{frame}

\begin{frame}{AES: Advanced Encryption Standard}
\textbf{Most widely used symmetric encryption (adopted 2001)}

\vspace{0.3cm}
\textbf{Specifications:}
\begin{itemize}
    \item Block size: 128 bits
    \item Key size: 128, 192, or 256 bits
    \item Rounds: 10, 12, or 14 (depending on key size)
    \item Based on substitution-permutation network
\end{itemize}

\vspace{0.3cm}
\textbf{Used everywhere:}
\begin{itemize}
    \item WiFi (WPA2/WPA3)
    \item VPNs and TLS/SSL
    \item File encryption
    \item Messaging apps
    \item Government TOP SECRET (with 256-bit keys)
\end{itemize}

% IMAGE SUGGESTION: AES encryption round diagram
\end{frame}

\begin{frame}{AES Security}
\textbf{AES-256 has $2^{256}$ possible keys}

\vspace{0.3cm}
Testing 1 trillion trillion keys per second:
\[
\frac{2^{256}}{10^{24}} \approx 3.7 \times 10^{51} \text{ seconds} \approx 10^{44} \text{ years}
\]

\vspace{0.3cm}
\begin{alertblock}{Key Insight}
Doubling key length doesn't double security - it \emph{squares} it!

$2^{256} = (2^{128})^2$ means 256-bit is $(2^{128})$ times more secure than 128-bit
\end{alertblock}
\end{frame}

% =====================================================
\section{Asymmetric Encryption}
% =====================================================

\begin{frame}{The Key Distribution Problem}
\textbf{Symmetric encryption challenge:}

\begin{itemize}
    \item Alice and Bob need shared secret key
    \item Never met in person
    \item Any key sent over internet could be intercepted
    \item Can't encrypt the key (need another key - infinite regress!)
\end{itemize}

\vspace{0.5cm}
\textbf{This problem plagued cryptography for centuries...}

% IMAGE SUGGESTION: Diagram showing Alice, Bob, and Eve with intercepted key
\end{frame}

\begin{frame}{The Revolutionary Idea}
\textbf{Public-Key Cryptography (Diffie-Hellman, 1976)}

\vspace{0.3cm}
\textbf{Core concept:}
\begin{itemize}
    \item Two mathematically related keys: \textbf{public} and \textbf{private}
    \item Public key shared freely
    \item Private key kept secret
    \item Encrypted with public → decrypted with private
    \item Computationally infeasible to derive private from public
\end{itemize}

\vspace{0.3cm}
\begin{block}{Mailbox Analogy}
Public key = mailbox (anyone can drop letters)\\
Private key = mailbox key (only you can open)
\end{block}

% IMAGE SUGGESTION: Mailbox analogy illustration
\end{frame}

\begin{frame}{How It Solves Key Distribution}
\begin{enumerate}
    \item Bob generates public-private key pair
    \item Bob publishes public key
    \item Alice looks up Bob's public key
    \item Alice encrypts message with Bob's public key
    \item Alice sends ciphertext
    \item Bob decrypts with his private key
\end{enumerate}

\vspace{0.3cm}
\textbf{No secret information exchanged beforehand!}

% IMAGE SUGGESTION: Step-by-step diagram of public key encryption process
\end{frame}

\begin{frame}{Symmetric vs. Asymmetric}
\begin{center}
\begin{tabular}{|l|l|l|}
\hline
\textbf{Property} & \textbf{Symmetric} & \textbf{Asymmetric} \\
\hline
Keys & Same key & Public + Private \\
Key distribution & Difficult & Easy \\
Speed & Very fast & Much slower \\
Key size & 128-256 bits & 2048-4096 bits \\
Use & Bulk encryption & Key exchange \\
Examples & AES & RSA, ECC \\
\hline
\end{tabular}
\end{center}

\vspace{0.3cm}
\textbf{In practice:} Combine both for best results!
\end{frame}

\begin{frame}{RSA Encryption}
\textbf{RSA (Rivest-Shamir-Adleman, 1977)}

Based on difficulty of factoring large numbers

\vspace{0.3cm}
\textbf{Key generation:}
\begin{enumerate}
    \item Choose two large primes $p$ and $q$
    \item Compute $n = p \times q$ (public)
    \item Compute $\phi(n) = (p-1)(q-1)$
    \item Choose public exponent $e$ (commonly 65537)
    \item Compute private exponent $d$
    \item Public key: $(n, e)$, Private key: $(n, d)$
\end{enumerate}

\vspace{0.3cm}
\textbf{Encryption:} $\text{Ciphertext} = \text{Message}^e \mod n$

\textbf{Decryption:} $\text{Message} = \text{Ciphertext}^d \mod n$

% IMAGE SUGGESTION: RSA key generation flowchart
\end{frame}

\begin{frame}{RSA in Practice}
\textbf{Key sizes:}
\begin{itemize}
    \item 1024 bits: \alert{Deprecated} (factored by nation-states)
    \item 2048 bits: Current minimum
    \item 3072 bits: High security
    \item 4096 bits: Very high security (slower)
\end{itemize}

\vspace{0.3cm}
\textbf{Common uses:}
\begin{itemize}
    \item TLS/SSL certificates (HTTPS)
    \item Email encryption (PGP/GPG)
    \item SSH keys
    \item Digital signatures
\end{itemize}

\vspace{0.3cm}
\textbf{Note:} Too slow for large data - use hybrid encryption!
\end{frame}

\begin{frame}{Elliptic Curve Cryptography (ECC)}
\textbf{Same security as RSA with much smaller keys}

\vspace{0.3cm}
\begin{center}
\begin{tabular}{|c|c|}
\hline
\textbf{RSA Key Size} & \textbf{ECC Key Size} \\
\hline
1024 bits & 160 bits \\
2048 bits & 224 bits \\
3072 bits & 256 bits \\
15360 bits & 512 bits \\
\hline
\end{tabular}
\end{center}

\vspace{0.3cm}
\textbf{Advantages:}
\begin{itemize}
    \item Smaller keys, certificates, signatures
    \item Faster operations
    \item Less bandwidth and storage
\end{itemize}

% IMAGE SUGGESTION: Elliptic curve graph
\end{frame}

\begin{frame}{ECC Applications}
\textbf{Where ECC is used:}
\begin{itemize}
    \item Bitcoin and cryptocurrency wallets
    \item Modern TLS/SSL certificates
    \item Signal, WhatsApp (Curve25519)
    \item Apple iMessage
    \item SSH keys
    \item IoT devices (limited computing power)
\end{itemize}

\vspace{0.3cm}
\textbf{Popular curves:}
\begin{itemize}
    \item P-256, P-384, P-521 (NIST standard)
    \item Curve25519 (modern, secure)
    \item Ed25519 (optimized for signatures)
\end{itemize}
\end{frame}

% =====================================================
\section{Cryptographic Hash Functions}
% =====================================================

\begin{frame}[fragile]{What is a Hash Function?}
\textbf{Takes any input → produces fixed-size output}

\vspace{0.3cm}
\begin{exampleblock}{Example: SHA-256}
\texttt{Input: "Hello, World!"}\\
\texttt{Hash: dffd6021bb2bd5b0af676290809ec3a5...}

\vspace{0.2cm}
\texttt{Input: "Hello, World?" (note the ?)}\\
\texttt{Hash: a0e0d1d23f58c99b8a73fdecd1ba0d8f...}
\end{exampleblock}

\vspace{0.3cm}
\textbf{Notice:}
\begin{itemize}
    \item Same length output (256 bits)
    \item Tiny change → completely different hash
    \item Looks random
\end{itemize}

% IMAGE SUGGESTION: Visual diagram of hash function with avalanche effect
\end{frame}

\begin{frame}{Properties of Cryptographic Hashes}
\begin{enumerate}
    \item \textbf{Deterministic:} Same input → same hash
    \item \textbf{Fast to compute}
    \item \textbf{Pre-image resistance:} Hard to find input from hash
    \item \textbf{Second pre-image resistance:} Hard to find different input with same hash
    \item \textbf{Collision resistance:} Hard to find any two inputs with same hash
    \item \textbf{Avalanche effect:} Small change → drastically different hash
    \item \textbf{Fixed output size}
\end{enumerate}
\end{frame}

\begin{frame}{Common Hash Functions}
\begin{center}
\begin{tabular}{|l|c|c|}
\hline
\textbf{Algorithm} & \textbf{Output Size} & \textbf{Status} \\
\hline
MD5 & 128 bits & \alert{Broken} \\
SHA-1 & 160 bits & \alert{Deprecated} \\
SHA-256 & 256 bits & \textcolor{green}{Secure} \\
SHA-512 & 512 bits & \textcolor{green}{Secure} \\
SHA-3 & Various & \textcolor{green}{Secure} \\
BLAKE2/3 & Various & \textcolor{green}{Secure} \\
\hline
\end{tabular}
\end{center}

\vspace{0.3cm}
\begin{alertblock}{Important}
Never use MD5 or SHA-1 for security purposes!
\end{alertblock}
\end{frame}

\begin{frame}{Hash Function Applications}
\textbf{1. Data Integrity / Checksums}
\begin{itemize}
    \item Verify file downloads
    \item Detect corruption or tampering
\end{itemize}

\vspace{0.3cm}
\textbf{2. Password Storage}
\begin{itemize}
    \item Never store plaintext passwords!
    \item Store hash of password
    \item Add salt to prevent rainbow tables
    \item Use specialized functions: Argon2, bcrypt
\end{itemize}

\vspace{0.3cm}
\textbf{3. Digital Signatures}
\begin{itemize}
    \item Hash document, then sign hash
\end{itemize}

% IMAGE SUGGESTION: Diagram showing password hashing with salt
\end{frame}

\begin{frame}[fragile]{Password Hashing with Salt}
\textbf{Salt = random string added before hashing}

\vspace{0.3cm}
\begin{exampleblock}{Example}
User 1: "password123" + Salt "a8f9e2" → Hash 1\\
User 2: "password123" + Salt "x3k8m1" → Hash 2
\end{exampleblock}

\vspace{0.3cm}
\textbf{Stored in database:}
\begin{center}
\begin{tabular}{|l|l|l|}
\hline
\textbf{User} & \textbf{Salt} & \textbf{Hash} \\
\hline
alice & a8f9e2 & 3f7a8c9d... \\
bob & x3k8m1 & 9e2b5f1c... \\
\hline
\end{tabular}
\end{center}

\textbf{Prevents:}
\begin{itemize}
    \item Rainbow table attacks
    \item Identifying users with same password
    \item Parallel attacks
\end{itemize}
\end{frame}

\begin{frame}{Blockchain and Cryptocurrencies}
\textbf{Hash functions are fundamental to blockchain:}

\vspace{0.3cm}
\begin{itemize}
    \item \textbf{Block hashing:} Each block contains hash of previous (creates chain)
    \item \textbf{Proof of work:} Find nonce where block hash starts with zeros
    \item \textbf{Addresses:} Bitcoin addresses are hashes of public keys
    \item \textbf{Merkle trees:} Efficiently verify transactions
\end{itemize}

% IMAGE SUGGESTION: Blockchain diagram showing hash links between blocks
\end{frame}

% =====================================================
\section{Digital Signatures}
% =====================================================

\begin{frame}{Digital Signatures}
\textbf{Electronic equivalent of handwritten signature}

\vspace{0.3cm}
\textbf{Provides:}
\begin{itemize}
    \item \textbf{Authentication:} Proves who sent message
    \item \textbf{Integrity:} Proves message not modified
    \item \textbf{Non-repudiation:} Sender cannot deny sending
\end{itemize}

\vspace{0.3cm}
\textbf{How it works:}
\begin{enumerate}
    \item Hash the message
    \item Encrypt hash with \textbf{private key} (signature)
    \item Send message + signature
    \item Receiver decrypts signature with \textbf{public key}
    \item Compare with hash of received message
\end{enumerate}

% IMAGE SUGGESTION: Digital signature process diagram
\end{frame}

\begin{frame}{Encryption vs. Signing}
\begin{center}
\begin{tabular}{|l|l|l|}
\hline
\textbf{Property} & \textbf{Encryption} & \textbf{Signing} \\
\hline
Lock with & Public key & Private key \\
Unlock with & Private key & Public key \\
Purpose & Confidentiality & Authentication \\
Who can do it? & Anyone & Only key owner \\
Who can verify? & Only key owner & Anyone \\
\hline
\end{tabular}
\end{center}

\vspace{0.5cm}
\textbf{Remember:}
\begin{itemize}
    \item \textbf{Encryption:} Public lock, private unlock (secrecy)
    \item \textbf{Signing:} Private lock, public unlock (authenticity)
\end{itemize}
\end{frame}

\begin{frame}{Public Key Infrastructure (PKI)}
\textbf{Problem:} How do you know a public key belongs to who it claims?

\vspace{0.3cm}
\textbf{Solution: Digital Certificates}

\vspace{0.3cm}
\textbf{Certificate contains:}
\begin{itemize}
    \item Public key
    \item Owner's identity
    \item Validity period
    \item Digital signature from Certificate Authority (CA)
\end{itemize}

\vspace{0.3cm}
\textbf{Trust chain:}
Website → Intermediate CA → Root CA (pre-installed in browser)

% IMAGE SUGGESTION: PKI certificate chain diagram
\end{frame}

\begin{frame}{How HTTPS Works}
\textbf{When you visit \texttt{https://www.example.com}:}

\begin{enumerate}
    \item Server sends certificate (with public key)
    \item Browser verifies:
    \begin{itemize}
        \item Certificate signed by trusted CA
        \item Not expired
        \item Matches domain name
        \item Not revoked
    \end{itemize}
    \item Browser and server perform key exchange
    \item Establish symmetric encryption key
    \item All data encrypted with session key
    \item Green padlock appears!
\end{enumerate}

\vspace{0.3cm}
\textbf{Provides:} Confidentiality + Authentication + Integrity

% IMAGE SUGGESTION: HTTPS handshake diagram
\end{frame}

% =====================================================
\section{Modern Applications}
% =====================================================

\begin{frame}{End-to-End Encryption (E2EE)}
\textbf{Only sender and recipient can read messages}

Not even the service provider!

\vspace{0.3cm}
\textbf{Examples:}
\begin{itemize}
    \item Signal
    \item WhatsApp
    \item iMessage
    \item ProtonMail
\end{itemize}

\vspace{0.3cm}
\textbf{Contrast with:}
\begin{itemize}
    \item Regular email: Provider can read
    \item SMS: Carrier can read
    \item Social media DMs: Platform can read
\end{itemize}

% IMAGE SUGGESTION: E2EE vs non-E2EE comparison diagram
\end{frame}

\begin{frame}{Password Managers}
\textbf{Securely store all passwords encrypted}

\vspace{0.3cm}
\textbf{How they work:}
\begin{itemize}
    \item Master password derives encryption key
    \item All passwords encrypted with AES-256
    \item Only you know master password
\end{itemize}

\vspace{0.3cm}
\textbf{Why use one:}
\begin{itemize}
    \item Unique strong password for every site
    \item Protection against phishing
    \item Protection against keyloggers
    \item No need to remember passwords
\end{itemize}

\vspace{0.3cm}
\textbf{Examples:} 1Password, Bitwarden, KeePass
\end{frame}

\begin{frame}{VPNs (Virtual Private Networks)}
\textbf{Create encrypted tunnels for internet traffic}

\vspace{0.3cm}
\textbf{What VPNs do:}
\begin{itemize}
    \item Encrypt traffic between you and VPN server
    \item Hide your IP address
    \item Bypass geographic restrictions
    \item Protect on public WiFi
\end{itemize}

\vspace{0.3cm}
\textbf{What VPNs don't do:}
\begin{itemize}
    \item Make you completely anonymous
    \item Protect against malware
    \item Encrypt data after VPN server
\end{itemize}

% IMAGE SUGGESTION: VPN tunnel diagram
\end{frame}

% =====================================================
\section{Threats and Future}
% =====================================================

\begin{frame}{Common Attacks}
\begin{description}
    \item[Brute Force] Try every possible key
    \item[Dictionary Attack] Try common passwords
    \item[Rainbow Tables] Precomputed password hashes
    \item[Man-in-the-Middle] Intercept communication
    \item[Side-Channel] Extract secrets via timing, power
    \item[Implementation Bugs] Heartbleed, POODLE, etc.
\end{description}

\vspace{0.3cm}
\textbf{Defense:} Large keys, salting, authentication, constant-time algorithms, updates
\end{frame}

\begin{frame}{The Quantum Computing Threat}
\textbf{Quantum computers could break current public-key crypto}

\vspace{0.3cm}
\textbf{Shor's Algorithm (1994):}
\begin{itemize}
    \item Can factor large numbers efficiently
    \item Breaks: RSA, Diffie-Hellman, ECC
\end{itemize}

\vspace{0.3cm}
\textbf{Symmetric encryption less affected:}
\begin{itemize}
    \item Grover's algorithm only quadratic speedup
    \item Double key size restores security
\end{itemize}

\vspace{0.3cm}
\textbf{Timeline:} 10-30+ years before threat is real

\textbf{Solution:} Post-Quantum Cryptography (PQC)

% IMAGE SUGGESTION: Timeline showing quantum computing development
\end{frame}

\begin{frame}{Post-Quantum Cryptography}
\textbf{NIST standardized quantum-resistant algorithms:}

\vspace{0.3cm}
\begin{itemize}
    \item \textbf{CRYSTALS-Kyber} (encryption)
    \item \textbf{CRYSTALS-Dilithium} (signatures)
    \item \textbf{FALCON} (signatures)
    \item \textbf{SPHINCS+} (signatures)
\end{itemize}

\vspace{0.3cm}
\textbf{Based on problems quantum computers can't solve:}
\begin{itemize}
    \item Lattice problems
    \item Hash-based schemes
    \item Multivariate polynomials
\end{itemize}
\end{frame}

\begin{frame}{Social Engineering}
\begin{alertblock}{Remember}
``The weakest link is often the human, not the cryptography.''
\end{alertblock}

\vspace{0.3cm}
\textbf{Common tactics:}
\begin{itemize}
    \item Phishing emails
    \item Pretexting (impersonation)
    \item Baiting (infected USB drives)
    \item Tailgating (physical access)
\end{itemize}

\vspace{0.3cm}
\textbf{Defense:}
\begin{itemize}
    \item Security awareness training
    \item Two-factor authentication
    \item Principle of least privilege
\end{itemize}

% IMAGE SUGGESTION: Phishing example or social engineering attack diagram
\end{frame}

% =====================================================
\section{Best Practices}
% =====================================================

\begin{frame}{Best Practices for Users}
\begin{enumerate}
    \item Use strong, unique passwords (12+ characters)
    \item Use a password manager
    \item Enable two-factor authentication (2FA)
    \item Keep software updated
    \item Use HTTPS (look for padlock)
    \item Be wary of phishing
    \item Use end-to-end encrypted messaging
    \item Encrypt sensitive data (full disk encryption)
    \item Use VPN on public WiFi
\end{enumerate}
\end{frame}

\begin{frame}{Best Practices for Developers}
\begin{enumerate}
    \item \textbf{Don't roll your own crypto!}
    \item Use established libraries (NaCl, OpenSSL)
    \item Use modern algorithms (AES-256, SHA-256, Argon2)
    \item Never store plaintext passwords
    \item Use secure random number generators
    \item Implement rate limiting
    \item Use TLS 1.3 with strong ciphers
    \item Keep dependencies updated
    \item Perform security audits
\end{enumerate}
\end{frame}

\begin{frame}{Common Mistakes to Avoid}
\begin{itemize}
    \item Using deprecated algorithms (MD5, SHA-1, DES)
    \item Storing passwords in plaintext
    \item Reusing keys, IVs, or nonces
    \item Using ECB mode
    \item Implementing your own crypto
    \item Trusting user input
    \item Ignoring timing attacks
    \item Hardcoding keys in source code
    \item Using weak random generators
\end{itemize}
\end{frame}

% =====================================================
\section{Conclusion}
% =====================================================

\begin{frame}{Key Takeaways}
\begin{enumerate}
    \item \textbf{Cryptography is essential} for our digital lives
    \item \textbf{Security is hard} - small mistakes = big consequences
    \item \textbf{Use established tools} - don't implement your own
    \item \textbf{Defense in depth} - crypto is necessary but not sufficient
    \item \textbf{Stay informed} - threats evolve constantly
    \item \textbf{Human element matters} - strongest crypto can't stop phishing
\end{enumerate}
\end{frame}

\begin{frame}{The Journey}
\textbf{From ancient ciphers to modern systems:}

\begin{itemize}
    \item Classical: Caesar, Vigenère, One-Time Pad
    \item Symmetric: AES, ChaCha20
    \item Asymmetric: RSA, ECC, Diffie-Hellman
    \item Hash Functions: SHA-256, Argon2
    \item Applications: HTTPS, E2EE, blockchain
    \item Future: Post-quantum cryptography
\end{itemize}

\vspace{0.5cm}
\textbf{Cryptography continues to evolve to meet new challenges!}
\end{frame}

\begin{frame}{Further Resources}
\textbf{Books:}
\begin{itemize}
    \item \emph{The Code Book} by Simon Singh
    \item \emph{Applied Cryptography} by Bruce Schneier
    \item \emph{Cryptography Engineering} by Ferguson et al.
\end{itemize}

\vspace{0.3cm}
\textbf{Online:}
\begin{itemize}
    \item cryptohack.org - Learn through challenges
    \item cryptopals.com - Practical exercises
    \item cryptii.com - Experiment with ciphers
    \item Coursera: Stanford Cryptography course
\end{itemize}
\end{frame}

\begin{frame}
\begin{center}
\Huge Thank You!

\vspace{1cm}
\large Questions?

\vspace{1cm}
\normalsize
\textit{``With great power comes great responsibility.''}

\vspace{0.5cm}
Use cryptography to protect privacy and security, not to cause harm.
\end{center}
\end{frame}

\end{document}
